\documentclass{article}

\usepackage{amsmath}

\begin{document}

{\bf Worksheet \#2; date: 08/29/2018}

{\bf MATH 55 Discrete Mathematics}

\begin{enumerate}
\item {\em True / False?} The negation of ``all heroes wear cape'' is ``there exists a hero who does not wear cape''.

\item {\em True / False?} All trees that can walk will also eat you.

\item {\em (Rosen 1.4.36c)} Find a counterexample, if possible to these universally quantified statements, where the domain for all variables consists of all integers.
\[
\forall x(|x| < 0)
\]

\item What do we need to do to show the following statement?
\[
	\text{There exists a prime number that is larger that $10^6$.}
\]

\item {\em (Rosen 1.4.43)} Determine whether $\forall x (P(x) \to Q(x))$ and $\forall x\, P(x) \to \forall x\, Q(x)$ are logically equivalent. Justify your answer.

\item {\em (Rosen 1.5.23d)} Express the following mathematical statements using predicates, quantifiers, logical connectives, and mathematical operators.
\[
	\text{A negative real number does not have a square root that is a real number.}
\]

\item {\em (Rosen 1.5.49a)} Show that $\forall x\, P(x) \wedge \exists x\, Q(x)$ equivalent to $\forall x \exists y (P(x) \wedge Q(y))$, where all quantifiers have the same nonempty domain.

\item {\em (Rosen 1.6.10e)} What relevant conclusions can be drawn from the following set of premise?
\begin{align*}
& \text{All foords that are healthy to eat do not taste good.} \\
& \text{Tofu is healthy to eat.} \\
& \text{You only eat what tastes good.} \\
& \text{You do not eat tofu.} \\
& \text{Cheeseburgers are not healthy to eat.}
\end{align*}

\item {\em (Rosen 1.6.19)} Determine whether each of these arguments is valid. If it is not correct, what logical error occurs?
\begin{enumerate}
\item If $n$ is a real number such that $n > 1$, then $n^2 > 1$. Suppose that $n^2 > 1$. Then $n > 1$.
\item If $n$ is a real number with $n > 3$, then $n^2 > 9$. Suppose that $n^2 \le 9$. Then $n \le 3$.
\item If $n$ is a real number with $n > 2$, then $n^2 > 4$. Suppose that $n \le 2$. Then $n^2 \le 4$.
\end{enumerate}
\end{enumerate}

\end{document}